\documentclass[11pt]{article}

\usepackage[utf8]{inputenc}
\usepackage[T1]{fontenc}
\usepackage{lmodern}
\usepackage{geometry}
\usepackage{amsmath,amssymb,amsfonts,amsthm,mathtools}
\usepackage{bm}
\usepackage{enumitem}
\usepackage{microtype}
\usepackage{hyperref}
\usepackage{titlesec}
\usepackage{booktabs}
\usepackage{array}
\usepackage{setspace}
\usepackage{mathrsfs}
\usepackage{physics}

\geometry{margin=1in}
\hypersetup{colorlinks=true, linkcolor=black, urlcolor=blue, citecolor=black}
\setlist[enumerate]{topsep=0.4em,itemsep=0.35em}
\setlist[itemize]{topsep=0.3em,itemsep=0.25em}
\titleformat{\section}{\large\bfseries}{\thesection.}{0.5em}{}
\titleformat{\subsection}{\normalsize\bfseries}{\thesubsection.}{0.5em}{}
\setstretch{1.06}

\newcommand{\Aether}{\AE{}ther}
\newcommand{\Aflow}{\AE{}ther-flow}
\newcommand{\TheoryName}{\Aflow{} Interpretation of Relativity}
\newcommand{\TheoryProgram}{\Aether{} / \Aflow{} framework}
\newcommand{\CompletedMetric}{g^{\sharp}}

\newcommand{\CurrentBodyImageCore}{\mathfrak S^{\mathrm{disp,resp,cbimg}}}
\newcommand{\CanonicalImageBody}[1]{\mathcal U_{#1}^{\mathrm{disp,resp,cbimg}}}
\newcommand{\CanonicalImageShell}[1]{\mathcal Z_{#1}^{\mathrm{disp,resp,cbimg}}}
\newcommand{\CanonicalImageResponseLaw}[1]{\mathscr Q_{#1}^{\mathrm{disp,resp,cbimg}}}

\newtheorem{definition}{Definition}
\newtheorem{proposition}{Proposition}
\newtheorem{remark}{Remark}
\newtheorem{corollary}{Corollary}
\newtheorem{theorem}{Theorem}

\title{\textbf{The \TheoryName{}}\\[0.4em]
\large Primitive-Reservoir Observer-Reduction\\
\large Section-Centered Hidden-Response\\
\large Canonical Image Response / Exact-Shell\\
\large Next Benchmark Criterion}
\author{}
\date{}

\begin{document}

\maketitle

\begin{abstract}
The current active-primary theorem now proves that, once the canonical image body \(\CanonicalImageBody{\sigma}\) and the restricted projection / canonical-lift relation on that body are held fixed, the canonical image response law \(\CanonicalImageResponseLaw{\sigma}\) and the canonical exact-shell image \(\CanonicalImageShell{\sigma}\) are already uniquely forced at benchmark scope. Another same-output deeper-origin relay that merely preserves that same image body together with the same image response and shell package no longer counts as the honest default next benchmark-facing manuscript.

This note records the routing consequence of that gain. The next honest benchmark-facing move must now derive or justify the restricted projection / canonical-lift relation itself from still more primitive explicit substrate structure, or derive the full canonical current-body image core \(\CurrentBodyImageCore\) in one genuinely cleaner deeper package, or prove a still sharper theorem that makes one of those moves unavoidable while preserving the same benchmark one-metric package, the same ordinary matter route, and the same claim boundary. The one operative metric remains exactly \(\CompletedMetric\), there is no second operative metric, the low-energy matter law is not enlarged, and no stronger promotion language is licensed.
\end{abstract}

\section{Introduction}

The active derivational program of \TheoryName{} is no longer merely at canonical-image-body scope in the broad sense recorded by the prior theorem. The preceding canonical-image response / exact-shell necessity / uniqueness theorem already showed that once the image body itself and the fixed restricted projection / canonical-lift relation are in hand, the response-level and exact-shell constituents of the current image core carry no remaining benchmark-facing presentational freedom. That theorem therefore spent the response and shell options that had remained open below the canonical-image-body frontier.

The present note records the routing consequence of that gain. It does not reopen the already-screened larger substrate-body presentation, the larger projection-operator core, the full displacement-response substrate law, the already-spent canonical image body, or the now-spent image response / shell package as if any of those were again independently live. It records only which kind of next manuscript would now count as an honest benchmark-facing gain once those options are no longer the open burden.

\paragraph{Framework and claim boundary.}
\TheoryProgram{} denotes the broader \Aether{} / \Aflow{} framework built on the claims that the \Aether{} is the underlying four-dimensional substrate of reality, the \Aflow{} is its intrinsic ordered motion, observed three-dimensional space is the local experiential slice of that deeper substrate, S-time is the experienced order of change arising from matter, light, and the \Aflow{}, and observed expansion is the three-dimensional appearance of deeper four-dimensional ordered motion. Within that framework, \TheoryName{} denotes the exact-closure theory in which the effective gravitational dynamics are adopted to be exactly Einsteinian with universal matter coupling. In the active sequence, \emph{adoption} means use of that established relativistic dynamics without claiming substrate derivation, while \emph{derivation} is reserved for a first-principles recovery from explicit substrate variables.

The present note stays strictly inside that boundary. It records only which kind of next manuscript would now count as an honest benchmark-facing gain below the canonical-image response / exact-shell necessity / uniqueness frontier.

\section{Fixed Inputs}

\begin{definition}[Canonical-image response / exact-shell frontier inputs]
The criterion recorded here uses the following fixed inputs.
\begin{enumerate}
    \item The benchmark exact-closure package which fixes the public one-metric GR target of the repository.
    \item The benchmark gatekeeping layer including the recorded safeguard that blocks same-output deeper-origin relays from counting as new front-line progress once the live benchmark-relevant datum is unchanged.
    \item The earlier theorem showing that the canonical image body \(\CanonicalImageBody{\sigma}\) is already uniquely forced inside the current image-core frontier.
    \item The later theorem proving that, once that body and the restricted projection / canonical-lift relation are fixed, the image response law \(\CanonicalImageResponseLaw{\sigma}\) and the exact-shell image \(\CanonicalImageShell{\sigma}\) are already uniquely forced as well.
\end{enumerate}
\end{definition}

\begin{remark}
The present note does not replace those manuscripts. It records the routing consequence of reading them together under the active safeguard against recursive depth drift.
\end{remark}

\section{Why the Live Burden Now Sits Below Response / Shell Restatement}

\begin{proposition}[The live burden is renewed below canonical-image response / shell restatement]
At the current stage of the primitive-reservoir observer-reduction line, the canonical-image response / exact-shell necessity / uniqueness theorem renews the live benchmark-facing burden below any mere restatement of the same image response / shell package on the fixed canonical image body.
\end{proposition}

\begin{proof}
That theorem changes benchmark-facing status because it proves that the image response law and exact-shell image are already uniquely forced once the canonical image body and the restricted projection / canonical-lift relation are held fixed. Once that uniqueness statement is in hand, another manuscript that merely preserves the same response / shell package no longer removes any remaining benchmark-facing freedom. The live burden therefore no longer sits on the task of restating or relocating \(\CanonicalImageResponseLaw{\sigma}\) or \(\CanonicalImageShell{\sigma}\). It sits on deriving the still-unspent restricted projection / canonical-lift relation from deeper explicit substrate structure, or on deriving the full image core in a genuinely cleaner deeper package.
\end{proof}

\begin{definition}[Benchmark-neutral same-response / shell relay below the current frontier]
A \emph{benchmark-neutral same-response / shell relay below the current frontier} is a manuscript that preserves the same benchmark one-metric output, the same ordinary matter route, the same canonical image body \(\CanonicalImageBody{\sigma}\), the same canonical image response law \(\CanonicalImageResponseLaw{\sigma}\), the same canonical exact-shell image \(\CanonicalImageShell{\sigma}\), the same canonical current-body image core \(\CurrentBodyImageCore\), and the same downstream benchmark-facing consequences while merely relocating the unexplained substrate layer deeper, renaming the same image-level constituents, restoring the already-screened larger substrate-body presentation, or re-expanding back to the larger projection-operator core, the full displacement-response substrate law, or the larger pullback-body-operator, pullback-operator, pullback-quadratic, hidden-response, residual-normal-form, or pre-proto-section presentations without a new derivation, new forcing theorem, or comparably sharp new gain.
\end{definition}

\section{The Current Post-Response / Shell Criterion}

\begin{definition}[Admissible next benchmark-facing manuscript below the current frontier]
Below the current canonical-image response / exact-shell necessity / uniqueness frontier, a manuscript counts as an admissible next benchmark-facing move only if it does at least one of the following:
\begin{enumerate}
    \item derives or justifies the restricted projection / canonical-lift relation on the fixed current displacement body from still more primitive explicit substrate structure in a genuinely benchmark-facing way;
    \item derives or justifies the full canonical current-body image core \(\CurrentBodyImageCore\) in one cleaner deeper package if that unified derivation is genuinely available;
    \item proves a sharper necessity, uniqueness, or compatibility theorem that makes one of the previous two moves unavoidable while preserving the same benchmark one-metric package, the same ordinary matter route, and the same claim boundary.
\end{enumerate}
\end{definition}

\begin{theorem}[Next honest benchmark-facing task below the current frontier]
The next honest benchmark-facing manuscript below the current canonical-image response / exact-shell necessity / uniqueness frontier must derive or justify the restricted projection / canonical-lift relation from still more primitive explicit substrate structure, derive the full canonical current-body image core \(\CurrentBodyImageCore\) in a genuinely cleaner deeper package, or prove a still sharper theorem that makes one of those moves unavoidable. A manuscript that only repeats the same canonical image body together with the same image response law and exact-shell image through another deeper relay, restores the larger substrate-body presentation, re-expands the discussion back to the larger projection-operator core or the full displacement-response substrate law, or replays the larger already-spent presentations without such a new gain does not satisfy the criterion.
\end{theorem}

\begin{proof}
By Proposition 1, the live burden already lies below any mere restatement of the same response / shell package on the fixed image body. Therefore a subsequent manuscript must improve on the already-forced image-core data rather than merely accompany them. A deeper-origin manuscript can do so only if it changes benchmark-facing status by more than depth alone. If it simply reproduces the same image body, the same response law, and the same shell image together with the same downstream outputs, the routing safeguard classifies it as benchmark neutral. The remaining honest alternatives are therefore exactly those stated in the definition.
\end{proof}

\section{What the Next Move Should Not Be}

\begin{proposition}[Screened next moves below the current frontier]
Below the current canonical-image response / exact-shell necessity / uniqueness frontier, none of the following is the default next benchmark-facing move:
\begin{enumerate}
    \item another same-output deeper-origin relay that merely preserves the current canonical image body \(\CanonicalImageBody{\sigma}\), the current canonical image response law \(\CanonicalImageResponseLaw{\sigma}\), and the current canonical exact-shell image \(\CanonicalImageShell{\sigma}\), and therefore merely reproduces the same current image core \(\CurrentBodyImageCore\) together with the same downstream benchmark-facing consequences;
    \item another restoration of the full substrate-body presentation as if the already-screened larger body freedom were again independently live;
    \item another replay of the full projection-operator core, the full displacement-response substrate law, or the larger pullback-body-operator, pullback-operator, pullback-quadratic, hidden-response, residual-normal-form, or pre-proto-section presentations as if those already-spent larger presentations were again the honest current burden;
    \item a manuscript that introduces a second operative metric, enlarges the ordinary matter law, or uses stronger promotion language.
\end{enumerate}
\end{proposition}

\begin{proof}
Item (1) fails the preceding criterion because it leaves the renewed live burden unchanged. Item (2) likewise fails because it restores a larger presentation whose relevant freedom is already screened out. Item (3) fails because it re-expands to larger already-spent presentations rather than deriving the still-unspent projection / lift relation or the full image core. Item (4) violates the fixed claim boundary of the benchmark package and therefore cannot count as a disciplined next step on the active line.
\end{proof}

\begin{remark}
This screening statement is not a denial that those deeper or alternative manuscripts may matter strategically. It is a routing statement: they are not the default way to spend the next benchmark-facing move below the current frontier unless they do the benchmark-facing work named above.
\end{remark}

\section{Conclusion}

The positive result of this note is routing discipline below the canonical-image response / exact-shell necessity / uniqueness frontier. The present theorem remains the live benchmark-facing gain because it already proves that, once the canonical image body and the restricted projection / canonical-lift relation are fixed, the image response law and the exact-shell image are uniquely forced on the recorded route. The earlier canonical-image-body necessity / uniqueness theorem remains preserved main-line history because it first removed the body constituent as an independently live burden.

The scope remains narrow and honest. The benchmark package is unchanged: the one operative metric is still exactly \(\CompletedMetric\), the ordinary matter route is unchanged, there is no second operative metric, and the low-energy matter law is not enlarged. Nothing here upgrades adoption to derivation or licenses stronger promotion language.

The next open burden is therefore clear. The next benchmark-facing manuscript should derive or justify the restricted projection / canonical-lift relation from still more primitive explicit substrate structure, or derive the full current image core in a genuinely cleaner deeper package, or prove a still sharper theorem that makes one of those moves unavoidable. Another same-output deeper relay that merely preserves the same image body, the same image response law, and the same exact-shell image, another restoration of the full substrate-body presentation, or another replay of the full projection-operator core or the full displacement-response substrate law is not the clean next manuscript target at current scope.

\end{document}
